%! Carrying Out a Test for a Population Mean — Unit 7, Lesson 7.5 — Lesson Plan
\documentclass[10pt]{article}
\usepackage{apstats-article}
\usepackage{apstats-boxes}
\newcommand{\TallMath}[1]{$\displaystyle #1\rule[-1.4em]{0pt}{3.2em}$}
\newcommand{\UnitNumberName}{Unit 7: Inference for Quantitative Data: Means \quad}
\newcommand{\LessonNumberName}{Lesson 7.5: Carrying Out a Test for a Population Mean}

\begin{document}
\begin{center}
    {\Large\bfseries \CourseName: \SchoolYear \\
    {\small\bfseries \UnitNumberName \LessonNumberName}}
\end{center}

\begin{tcolorbox}[colback=sky, colframe=navy, boxrule=0.9pt, arc=2mm,
    left=2mm, right=2mm, top=1.5mm, bottom=1.5mm]
    \textbf{Primary Objective:} Students will carry out a one-sample $t$-test or
    matched-pairs $t$-test by computing the $t$-statistic, finding and interpreting the
    $p$-value, and making and justifying a formal conclusion about a population mean ---
    using the decision rule that compares $p$ to the significance level $\alpha$.
    (Big Ideas: VAR, DAT)
\end{tcolorbox}

\renewcommand{\arraystretch}{1.6}
\begin{skillbox}[Priority Ideas \& Skills]{goldbox}
\begin{minipage}[t]{0.48\linewidth}
    \textbf{AP Skill 3 --- Using Probability \& Simulation}
    \begin{itemize}
        \item Calculate the $t$-test statistic $t = \dfrac{\bar x - \mu_0}{s/\sqrt{n}}$
              and find the $p$-value (Skill 3.E; VAR-7.E).
    \end{itemize}
    \textbf{AP Skill 4 --- Statistical Argumentation}
    \begin{itemize}
        \item Interpret the $p$-value of a one-sample $t$-test in context, recognizing
              it is computed \emph{assuming $H_0$ is true} (Skill 4.B; DAT-3.E).
        \item Justify a claim about the population by comparing $p$ to $\alpha$ and
              writing a conclusion in context (Skill 4.E; DAT-3.F).
    \end{itemize}
\end{minipage}
\hfill
\begin{minipage}[t]{0.48\linewidth}
    \textbf{Key Understandings}
    \begin{itemize}
        \item The $t$-statistic measures how many standard errors $\bar x$ is from $\mu_0$,
              using $s/\sqrt{n}$ as the standard error.
        \item The $p$-value is the probability of a result as extreme or more extreme,
              \emph{given} $H_0$ is true --- not the probability $H_0$ is true.
        \item Decision rule: if $p \le \alpha$, reject $H_0$; if $p > \alpha$, fail to
              reject $H_0$. Never ``accept'' $H_0$.
        \item The conclusion restates the decision in terms of the original research question,
              with appropriate uncertainty language (``at the $\alpha$ significance level'').
    \end{itemize}
\end{minipage}
\end{skillbox}

\begin{skillbox}[{Vocabulary, Concepts, \& Theorems}]{greenbox}
\begin{minipage}[t]{0.48\linewidth}
    \begin{itemize}
        \item \textbf{$t$-statistic} --- $t = \dfrac{\bar x - \mu_0}{s/\sqrt{n}}$; measures
              distance of $\bar x$ from $\mu_0$ in standard error units
        \item \textbf{Standard error of $\bar x$} --- $SE_{\bar x} = s/\sqrt{n}$; estimates
              $\sigma_{\bar x}$ when $\sigma$ is unknown
        \item \textbf{$p$-value} --- probability of observing a $t$-statistic as extreme or
              more extreme than the computed one, assuming $H_0$ is true
    \end{itemize}
\end{minipage}
\hfill
\begin{minipage}[t]{0.48\linewidth}
    \begin{itemize}
        \item \textbf{Significance level ($\alpha$)} --- the threshold for rejection; typically
              0.05 unless otherwise stated
        \item \textbf{Reject $H_0$} --- if $p \le \alpha$; data provide sufficient evidence
              against $H_0$
        \item \textbf{Fail to reject $H_0$} --- if $p > \alpha$; data do not provide sufficient
              evidence against $H_0$ (this is NOT the same as proving $H_0$ true)
        \item \textbf{Four-step structure} --- State, Plan, Do, Conclude (SPDC)
    \end{itemize}
\end{minipage}
\end{skillbox}

\begin{skillbox}[Activate Prior Knowledge \& Spiral Review]{sky}
\begin{minipage}[t]{0.55\linewidth}\vspace{0pt}
    \textbf{Spiral Review (warm-up)}
    \begin{itemize}
        \item Lesson 7.4: identify the $t$-test, state $H_0$ and $H_a$, verify conditions
        \item Unit 6: the four-step inference structure (State, Plan, Do, Conclude) for a
              $z$-test --- same structure applies here
        \item Reading $t$-tables: find $p$-value from $t$-statistic and $df$ (or use
              \texttt{tcdf} on the TI-84)
    \end{itemize}
\end{minipage}
\hfill
\begin{minipage}[t]{0.38\linewidth}\vspace{0pt}\centering
    % Authored warm-up compiles to target/ --- no source PDF to embed here.
\end{minipage}
\end{skillbox}

\begin{skillbox}[Hook]{sky}
\begin{minipage}[t]{0.55\linewidth}
    \textbf{Finishing the Cereal Box Test (3--4 min)}\\
    Return to the hook from Lesson 7.4: a consumer group weighed 25 cereal boxes labeled
    18 oz and got $\bar x = 17.6$ oz, $s = 0.8$ oz. We already set up the test:
    $H_0: \mu = 18$, $H_a: \mu < 18$, conditions verified.\\[4pt]
    Now ask: ``What is the actual $t$-statistic? Is $\bar x = 17.6$ oz far enough below 18 oz
    to convince us the company is under-filling?''\\[4pt]
    Reveal: $t = \dfrac{17.6 - 18}{0.8/\sqrt{25}} = \dfrac{-0.4}{0.16} = -2.50$.\\[4pt]
    ``How do we turn this $t$ into a probability (a $p$-value)?''
\end{minipage}
\hfill
\begin{minipage}[t]{0.40\linewidth}
    \textbf{What You Already Know}
    \begin{itemize}
        \item The $t$-statistic formula is structurally identical to the $z$-statistic:
              $\text{statistic} = \dfrac{\text{estimate} - \text{null value}}{\text{SE}}$
        \item The $p$-value always answers: ``If $H_0$ were true, how likely is a result
              this extreme (or more so)?''
        \item The decision rule ($p \le \alpha$ $\Rightarrow$ reject) is the same as
              in Unit 6.
    \end{itemize}
\end{minipage}
\end{skillbox}

\begin{skillbox}[Lesson: Computing $t$ and Finding $p$]{sky}
\begin{minipage}[t]{0.48\linewidth}
    \textbf{The $t$-Statistic}
    \[
      t = \frac{\bar x - \mu_0}{s/\sqrt{n}} \quad \text{with } df = n - 1
    \]
    \begin{itemize}
        \item $\bar x - \mu_0$: how far the sample mean is from the null value
        \item $s/\sqrt{n}$: standard error of $\bar x$ (estimated, since $\sigma$ is unknown)
        \item The result follows a $t$-distribution with $df = n - 1$ when conditions are met
    \end{itemize}
    \textit{Cereal: $t = -2.50$, $df = 24$.}\\
    Interpretation: $\bar x = 17.6$ is 2.50 standard errors \emph{below} the null value of 18 oz.
\end{minipage}
\hfill
\begin{minipage}[t]{0.48\linewidth}
    \textbf{Finding the $p$-Value}
    \begin{itemize}
        \item \textbf{One-sided ($H_a: \mu < \mu_0$):} $p = P(t < t_{\text{obs}})$ with $df = n-1$
        \item \textbf{One-sided ($H_a: \mu > \mu_0$):} $p = P(t > t_{\text{obs}})$
        \item \textbf{Two-sided ($H_a: \mu \ne \mu_0$):} $p = 2 \cdot P(t > |t_{\text{obs}}|)$
    \end{itemize}
    \textbf{On TI-84:} \texttt{tcdf(lower, upper, df)}\\
    \textit{Cereal: $p = $ \texttt{tcdf}$(-10^{99}, -2.50, 24) \approx 0.010$.}\\[4pt]
    \textbf{Table B:} find $df = 24$ row; $t = 2.50$ falls between $t^* = 2.492$ ($p = 0.01$, one-tail).
\end{minipage}
\end{skillbox}

\begin{skillbox}[Lesson (cont.): Interpret and Conclude]{sky}
\begin{minipage}[t]{0.48\linewidth}
    \textbf{Interpreting the $p$-Value (DAT-3.E)}\\
    The $p$-value is the probability of observing a sample mean as far below 18 oz as 17.6 oz
    (or farther), \emph{if the true mean weight really is 18 oz}.\\[6pt]
    \textit{Cereal: $p \approx 0.010$.} ``If boxes truly average 18 oz, there's only a 1.0\%
    chance of seeing $\bar x \le 17.6$ oz in a sample of 25 by chance alone.''\\[4pt]
    Common error: ``There is a 1\% chance $H_0$ is true.'' This is wrong --- $p$ is computed
    \emph{assuming} $H_0$ is true.
\end{minipage}
\hfill
\begin{minipage}[t]{0.48\linewidth}
    \textbf{Formal Decision and Conclusion (DAT-3.F)}\\
    Compare $p$ to $\alpha = 0.05$:
    \begin{itemize}
        \item $p = 0.010 \le 0.05$ $\Rightarrow$ \textbf{Reject $H_0$}.
    \end{itemize}
    \textbf{Conclusion template:}\\
    ``Because $p = 0.010 < \alpha = 0.05$, we reject $H_0$. We have convincing evidence that
    the true mean weight of cereal boxes is less than 18 oz.''\\[4pt]
    \textit{Key phrases:} ``convincing evidence that [context of $H_a$]'' when rejecting;
    ``not sufficient evidence that [\ldots]'' when failing to reject.
\end{minipage}
\end{skillbox}

\begin{skillbox}[Active Monitoring]{redbox}
\begin{minipage}[t]{0.48\linewidth}
    \textbf{Circulate and Check}
    \begin{itemize}
        \item $t$-formula: students use $s$ (not $\sigma$) in the denominator.
        \item $p$-value direction matches $H_a$ (left-tail for $<$, right-tail for $>$,
              two-tail for $\ne$).
        \item $p$-value interpretation references $H_0$ being true, not the probability
              that $H_0$ is true.
        \item Conclusion is in context, references $\alpha$, and does not say ``accept $H_0$.''
        \item Matched pairs: $t$-statistic computed on $\bar d$ and $s_d$, not on original data.
    \end{itemize}
\end{minipage}
\hfill
\begin{minipage}[t]{0.48\linewidth}
    \textbf{Cold-Call Prompts}
    \begin{itemize}
        \item ``What does the denominator $s/\sqrt{n}$ represent?''
        \item ``The $p$-value is 0.07 and $\alpha = 0.05$. What do you conclude?''
        \item ``Can we say the null hypothesis is true if we fail to reject it? Why not?''
        \item ``If $H_a$ is two-sided, how does the $p$-value calculation change?''
    \end{itemize}
\end{minipage}
\end{skillbox}

\begin{skillbox}[Group Work \& Differentiation]{redbox}
\begin{multicols}{3}
    \textbf{Tier R --- Remediate}
    \begin{itemize}
        \item Provide a formula strip: $t = (\bar x - \mu_0)/(s/\sqrt{n})$ and a completed
              $p$-value example from the TI-84.
        \item Focus on plugging in values and writing the conclusion sentence.
        \item Scenarios A and B only.
    \end{itemize}
    \columnbreak
    \textbf{Tier A --- Approaching}
    \begin{itemize}
        \item Complete all three scenarios from hook to conclusion.
        \item Interpret the $p$-value in a sentence for each.
    \end{itemize}
    \columnbreak
    \textbf{Tier E --- Extension}
    \begin{itemize}
        \item Complete all scenarios including the matched-pairs problem.
        \item For one scenario: explain how a larger $n$ would change $SE$ and $t$,
              and predict whether the $p$-value would increase or decrease.
    \end{itemize}
\end{multicols}
\end{skillbox}

\begin{skillbox}[Individual Work \& Assessment]{redbox}
\begin{minipage}[t]{0.48\linewidth}
    \textbf{Exit Ticket (5 min)}\\
    Context: Mean commute time. $H_0: \mu = 30$ min, $H_a: \mu > 30$ min.
    Random sample of 22 commuters: $\bar x = 33.4$ min, $s = 6.8$ min; distribution is
    approximately normal.
    \begin{enumerate}
        \item Compute $t$.
        \item Find the $p$-value (assume $p \approx 0.022$).
        \item Write the conclusion at $\alpha = 0.05$.
    \end{enumerate}
\end{minipage}
\hfill
\begin{minipage}[t]{0.48\linewidth}
    \textbf{AP-Style Check}\\
    \textit{A student computes $t = -1.85$ with $df = 19$ for a two-sided test. The
    $p$-value is approximately 0.08. At $\alpha = 0.05$, the correct conclusion is:}
    \begin{enumerate}[label=(\Alph*)]
        \item Reject $H_0$; $p < \alpha$.
        \item Fail to reject $H_0$; there is not sufficient evidence.
        \item Accept $H_0$; the data support the null.
        \item Reject $H_a$; the data are consistent with $H_0$.
    \end{enumerate}
    Answer: (B)
\end{minipage}
\end{skillbox}

\begin{skillbox}[Reinforcement \& Extension]{goldbox}
\begin{minipage}[t]{0.48\linewidth}
    \textbf{Homework Overview}
    \begin{itemize}
        \item Three full SPDC problems: one one-sided, one two-sided, one matched-pairs.
        \item Students compute $t$, find $p$, write interpretation and conclusion.
        \item Reflection: ``Explain, in plain language, what it means to fail to reject
              $H_0$. Is this the same as proving the null is true?''
    \end{itemize}
\end{minipage}
\hfill
\begin{minipage}[t]{0.48\linewidth}
    \textbf{Extension / Preview}
    \begin{itemize}
        \item \textit{Extension:} For a two-sided test with $t = 2.10$ and $df = 15$,
              find the $p$-value using both Table B and the TI-84. Compare.
        \item \textit{Preview:} Lesson 7.6 examines Type I and Type II errors and the
              power of a significance test --- consequences of the decisions we made today.
    \end{itemize}
\end{minipage}
\end{skillbox}

\end{document}
